\section{\label{sec:concs}结论}

部分子分布函数（PDFs）借核子的组分夸克与胶子刻画核子结构。这些 PDFs 定义为光前波函数的矩阵元，无法直接在欧几里得格点 QCD 中计算。然而原则上，PDFs 的 Mellin 矩可以通过
扭度2算符的矩阵元在格点 QCD 中计算。
遗憾的是，这类计算仅限于头几个矩，因为
格点调节子的超立方对称性在不同质量量纲的扭度2算符之间诱导幂发散混合，模糊了格点上确定的矩阵元的连续极限。当前对 PDFs 的测定依赖于对众多实验道数据的全局分析，
尚缺乏从第一性原理出发的 PDF 测定。

最近，Ji 提出了在格点 QCD 中确定 PDFs 的新途径；随后，Qiu 与 Ma 通过一个相关框架也提出了类似方案。
在这一途径中，人们在大核子
动量处计算欧几里得准 PDFs。两项近期的格点计算给出了有希望的结果，
但该方法的若干方面仍有待充分理解。第一，存在与格点计算中有限核子动量相关的系统不确定度这一实际问题。
随着计算资源的改进与已在进行的算法进展，
这一问题有望得到解决，其精度将足以让结果能在实验数据不足之处为全局分析提供有用输入。第二，有待澄清的理论问题有二：
定义准 PDFs 的延展算符的重正化；
以及欧几里得准 PDF 与光前 PDF 之间的关系——迄今后者仅通过单圈微扰论中的因子化公式得到分析。

我们通过引入由梯度流涂抹的场构造的准 PDF，回应了上述第一个理论考量。
我们明确证明：一旦并入核子质量修正，
并且在流时间相对核子动量的倒数很小的条件下，
涂抹欧几里得准 PDF 的 Mellin 矩与
光前 PDF 的重正化 Mellin 矩之间存在简单关系。
对该关系的修正在 ${\cal O}(\Lambda_{\mathrm{QCD}}^2/P_z^2)$（$P_z$ 为核子的欧几里得动量）与 ${\cal O}(\sqrt{\tau} \Lambda_{\mathrm{QCD}})$（$\tau$ 为流时间）量级上出现。由这一对应可知，只要
$\Lambda_{QCD},M_N\ll P_z \ll \tau^{-1/2}$，准 PDF 与光前 PDF
便可通过卷积关系进行匹配。

我们方法的主要优点是：梯度流使准
PDF 在连续极限下有限，并回避了
定义格点上准 PDF 的非局域算符的重正化问题。由此得到的连续矩阵元
与开展格点 QCD 计算所选用的离散作用量无关，并可用
连续时空微扰论直接匹配到 $\ms$ 方案下的相应光前 PDFs。结合非微扰步标度（step-scaling）程序，
该匹配可在足够高的能量处进行，
使微扰截断误差不再失控。我们方案的非微扰实现正在进行之中。


